\documentclass[11pt]{article}

\usepackage[margin=1in]{geometry}
\usepackage{amsmath,amssymb,amsthm}
\usepackage{underscore}   % lets identifier_names appear in \texttt{} without escaping
\usepackage{hyperref}
\usepackage{enumitem}
\usepackage{mathpartir}

\newtheorem{theorem}{Theorem}[section]
\newtheorem{lemma}[theorem]{Lemma}
\newtheorem{corollary}[theorem]{Corollary}
\theoremstyle{definition}
\newtheorem{definition}[theorem]{Definition}
\theoremstyle{remark}
\newtheorem*{remark}{Remark}

\newcommand{\inN}{\!\in\!}
\newcommand{\lean}[1]{\hfill{\footnotesize\upshape\texttt{[#1]}}\par\noindent\ignorespaces}
\newcommand{\leanpart}[1]{{\unskip\nobreak\hfil\penalty50\hskip1em\hbox{}\nobreak\hfill
  \mbox{\footnotesize\upshape\texttt{[#1]}}}}
\newcommand{\leanform}[1]{\par\smallskip\noindent{\upshape\small\textit{Lean form.} #1}}
\newcommand{\RId}{\mathsf{RId}}
\newcommand{\OId}{\mathsf{OId}}
\newcommand{\Rgn}{\mathsf{Rgn}}
\newcommand{\Stk}{\mathsf{Stk}}
\newcommand{\Open}{\mathit{Open}}
\newcommand{\Closed}{\mathit{Closed}}
\newcommand{\loc}[2]{\mathrm{loc}_{#1}(#2)}
\newcommand{\refs}{\mathit{refs}}
\newcommand{\oids}{\mathit{objectIds}}
\newcommand{\rid}{\mathit{regionId}}
\newcommand{\omap}{\mathit{objMap}}
\newcommand{\vmap}{\mathit{varMap}}
\newcommand{\bridge}{\mathit{bridgeObjectId}}
\newcommand{\status}{\mathit{status}}
\newcommand{\dom}{\mathrm{dom}}
\newcommand{\cat}{\mathbin{+\!\!+}}
\newcommand{\perm}{\equiv}
\newcommand{\sublist}{\sqsubseteq}
\newcommand{\noloc}{\bot}
\newcommand{\inv}[1]{\textit{#1}}

\title{Preservation of \textit{ValidConfig}: human-readable proofs}
\author{Hoang Nguyen}
\date{}

\begin{document}
\maketitle

\begin{abstract}
Pen-and-paper versions of the proofs in \texttt{lean/gc/Gc/Model/Preservation/}: for each mutation
operation, that it maps a \inv{ValidConfig} to a \inv{ValidConfig}. Each proof follows the
structure of its Lean counterpart, whose name is given in brackets on the right.
\end{abstract}

\tableofcontents

%======================================================================
\section{Preliminaries}
%======================================================================

\subsection{Configurations}

A configuration is $\sigma = \langle S, H\rangle$. The stack $S = [F_0, \dots, F_{n-1}]$ is a list of
frames whose \emph{top is the end of the list}; we write $S :: F$ for the stack with $F$ pushed on top
of $S$, i.e.\ the list $S \cat [F]$. The \emph{index} of a frame is its position, so $F_i$ has index
$i$ and a larger index is a later frame. The Lean \texttt{stackWithIndex} is the list of pairs
$(F_i, i)$; we write $(G, i) \in S$ for ``$i < |S|$ and $S[i] = G$''. In $S :: F$ the top frame $F$
has index $|S|$, and every frame of $S$ keeps the index it has in $S$.

The heap $H$ is a finite map from region ids to regions; a region $R$ has a $R.\bridge$, an
$R.\omap$ and a $R.\status \in \{\Open, \Closed\}$. A frame $G$ has a $G.\rid$, a bridge variable,
a $G.\omap$ of temporaries and a $G.\vmap$. In Lean, maps are association lists and the derived lists
below are concatenations in list order. Every invariant either asks for a list to be duplicate-free or
counts or tests membership in it, so all it sees of a list is its multiset of elements. We write
$\ell \perm \ell'$ for ``$\ell$ is a permutation of $\ell'$'' and $\ell \sublist \ell'$ for ``$\ell$
is a sublist of $\ell'$''; duplicate-freeness, counts and membership are invariant under $\perm$, and
are inherited (duplicate-freeness, membership) or bounded (counts) along $\sublist$.

\begin{definition}[References and object ids]
\lean{Gc/Model/Helpers.lean: Object.refs, \dots, RuntimeConfig.refs; Frame.objectIds, \dots,
RuntimeConfig.objectIds}
\begin{itemize}[nosep]
  \item $\refs(o)$ of an object $o$: its field values. $\refs(R)$ of a region: the $\refs$ of all its
    objects. $\refs(G)$ of a frame: the $\refs$ of all objects in $G.\omap$, followed by the values
    of $G.\vmap$.
  \item $\refs(S)$, $\refs(H)$: concatenation over frames, resp.\ regions.
    $\refs(\sigma) = \refs(S) \cat \refs(H)$.
  \item $\oids(G) = \dom(G.\omap)$, $\oids(R) = \dom(R.\omap)$; $\oids(S)$, $\oids(H)$ by
    concatenation; $\oids(\sigma) = \oids(S) \cat \oids(H)$.
\end{itemize}
Neither the bridge variable of a frame nor the $\bridge$ of a region counts as a reference.
\end{definition}

\begin{definition}[Location]
\label{def:loc}
\lean{Gc/Model/Helpers.lean: Reference.loc?}
For a reference $r$, $\loc{\sigma}{r}$ is a location $\Rgn(rid)$ or $\Stk(i)$, or undefined
($\noloc$):
\begin{itemize}[nosep]
  \item $\loc{\sigma}{\RId(rid)} = \Rgn(rid)$ if $rid \in \dom(H)$, and $\noloc$ otherwise.
  \item For $\loc{\sigma}{\OId(oid)}$, two searches run independently. The stack search looks from
    the top for the \emph{last} frame $(G,i)$ with $oid \in \dom(G.\omap)$. The heap search looks in
    entry order for the \emph{first} region $(rid, R)$ with $oid \in \dom(R.\omap)$. If only the stack search succeeds the location is
    $\Stk(i)$; if only the heap search succeeds it is $\Rgn(rid)$; if both or neither succeed it is
    $\noloc$.
\end{itemize}
\end{definition}

The ten invariants \inv{L1}, \inv{L2}, \inv{H1}, \inv{H2}, \inv{H3}, \inv{S1}, \inv{S2},
\inv{S3}, \inv{HS1}, \inv{HS2} are as in \texttt{main.tex}; $\sigma \vdash \inv{ValidConfig}$ is
their conjunction. The two we use most are:
\begin{itemize}[nosep]
  \item \inv{L1}: $\oids(\sigma)$ is duplicate-free;
  \item \inv{S1}: the list $[G.\rid \mid G \in S]$ is duplicate-free.
\end{itemize}

\subsection{Shared lemmas}

Each part of a lemma below is tagged with the Lean theorem that proves it. Where the Lean statement is
more general than what we need, or phrased differently, a \textit{Lean form} note gives the exact
shape.

\begin{lemma}[Locations under \inv{L1}]
\label{lem:loc-l1}
\lean{Gc/Model/Validity.lean}
If $\sigma \vdash \inv{L1}$ then for every $oid$:
\begin{enumerate}[nosep,label=(\alph*)]
  \item $\loc{\sigma}{\OId(oid)} = \Rgn(rid)$ iff $H(rid) = R$ for some $R$ with
    $oid \in \dom(R.\omap)$; \leanpart{oid\_loc\_rgn\_iff\_in\_heap}
  \item $\loc{\sigma}{\OId(oid)} = \Stk(i)$ iff $(G, i) \in S$ for some $G$ with
    $oid \in \dom(G.\omap)$. \leanpart{oid\_loc\_stk\_iff\_in\_stack}
\end{enumerate}
\leanform{Both take $\mathit{ValidConfig}\ \sigma$ as hypothesis but use only its \inv{L1} field.
In (a), ``$H(rid) = R$'' is \texttt{cfg.heap.lookup rid = some R}; in (b), ``$(G,i) \in S$'' is
\texttt{cfg.stackWithIndex[i]? = some G}.}
\end{lemma}
\begin{proof}
By \inv{L1}, $oid$ occurs at most once in $\oids(S) \cat \oids(H)$, so it lies in the $\omap$ of at
most one frame or region in total. The ``only if'' directions just read off the search that succeeded
in Definition~\ref{def:loc}. For ``if'' in (a): the heap search must find $(rid, R)$, because no other
region contains $oid$, and the stack search fails, because no frame contains it either; so the
location is $\Rgn(rid)$. Part (b) is symmetric: the stack search finds $(G,i)$ and the heap search
fails.
\end{proof}

\begin{lemma}[Status update]
\label{lem:status-update}
\lean{Gc/Model/Preservation/Common.lean}
Let $H(rid_0) = R_0$, let $s$ be a status and $H' = H[rid_0 \mapsto R_0']$ where $R_0'$ is $R_0$ with
its status set to $s$. Then:
\begin{enumerate}[nosep,label=(\alph*)]
  \item $\dom(H') = \dom(H)$; \leanpart{heap\_insert\_keys\_mem}
  \item $\refs(H') \perm \refs(H)$ \leanpart{heap\_insert\_refs\_perm}\\
    and $\oids(H') \perm \oids(H)$; \leanpart{heap\_insert\_objectIds\_perm}
  \item for every $rid$, $H'(rid)$ is defined iff $H(rid)$ is, and the two regions have the same
    $\omap$ (so the same $\refs$ and $\oids$); \leanpart{heap\_status\_update\_entries\_mem\_iff}
  \item if $\langle S, H\rangle \vdash \inv{L1}$, then for every $oid$ the heap search of
    Definition~\ref{def:loc} succeeds in $H'$ iff it succeeds in $H$, and then returns the same
    region id. \leanpart{heap\_status\_update\_find\_fst}
\end{enumerate}
\leanform{(a) and (b) hold for an insert of \emph{any} region at a present key, not just a status
change: (a) needs nothing more, (b) needs the new region to have the same $\refs$ (resp.\ a
permutation of the old $\omap$ keys), which a status change satisfies with equality. Other
operations reuse them in that generality. (c) is stated over heap entries: for all $r$ and $m$,
$\exists Q.\ (r, Q) \in H'.\mathit{entries} \wedge Q.\omap = m$ iff the same holds in $H$. (d) is
stated as an equation on the search results projected to their keys:
$(H'.\mathit{entries}.\mathtt{find?}\ p).\mathtt{map}\ \mathtt{fst} =
(H.\mathit{entries}.\mathtt{find?}\ p).\mathtt{map}\ \mathtt{fst}$, where $p$ tests
$oid \in \dom(Q.\omap)$.}
\end{lemma}
\begin{proof}
(a) $rid_0$ is already a key of $H$. (c) At $rid_0$ the region changes only in its status; at every
other key it does not change. (b) Inserting at a present key moves that entry to the front and
replaces its value with one of the same $\omap$; the derived lists are therefore rearranged but have
the same elements. (d) By (c) the regions of $H'$ containing $oid$ sit at exactly the keys of the
regions of $H$ containing $oid$. By \inv{L1} there is at most one such key, so the first-match search
returns it in both heaps, even though the entry order differs.
\end{proof}

%======================================================================
\section{\texttt{exit}}
%======================================================================

\begin{mathpar}
  \inferrule* [lab=exit]
    {|S| \ge 1 \\ H(F.\rid) = R \\ R.\status = \Open}
    {\langle S :: F, H \rangle \xrightarrow{\texttt{exit}} \langle S,\ H[F.\rid \mapsto R'] \rangle
     \quad\text{where } R' = R \text{ with } \status := \Closed}
\end{mathpar}

\texttt{exit} pops the top frame $F$ and closes its region. Throughout this section fix
$\sigma = \langle S :: F, H\rangle$ with $\sigma \vdash \inv{ValidConfig}$, and
$\sigma' = \langle S, H'\rangle$ with $H' = H[F.\rid \mapsto R']$ as above. Let $n = |S|$, the
index of $F$ in $\sigma$. This is exactly what \texttt{exit\_cases} extracts from
$\texttt{exit}\ \sigma = \mathsf{some}\ \sigma'$. Lemma~\ref{lem:status-update} applies with
$rid_0 = F.\rid$ and $s = \Closed$.

The proof has three parts: one lemma that describes every location of $\sigma'$ in terms of $\sigma$
(Lemma~\ref{lem:exit-loc}), the observation that nothing new is stored (Lemma~\ref{lem:exit-refs}),
and the ten invariants.

\subsection{Supporting lemmas}

\begin{lemma}[Locations after \texttt{exit}]
\label{lem:exit-loc}
\lean{exit\_loc}
For every reference $r$, either $\loc{\sigma'}{r} = \loc{\sigma}{r}$, or $\loc{\sigma'}{r} = \noloc$
and $\loc{\sigma}{r} = \Stk(i)$ for some $i$.
\end{lemma}
\begin{proof}
For $r = \RId(rid)$: the location depends only on $\dom$, which is unchanged by
Lemma~\ref{lem:status-update}(a), so $\loc{\sigma'}{r} = \loc{\sigma}{r}$.

For $r = \OId(oid)$: by Lemma~\ref{lem:status-update}(d), the heap search gives the same answer in
$H'$ as in $H$, so only the stack search can differ. The frames of $\sigma'$ are those of $\sigma$
minus the top one, with the same indices. Searching $\sigma$'s stack from the top therefore first
looks at $(F, n)$, and if $oid \notin \dom(F.\omap)$ then continues exactly as the search of
$\sigma'$'s stack.
\begin{itemize}[nosep]
  \item If $oid \notin \dom(F.\omap)$, both searches return the same result, so the two locations
    agree.
  \item If $oid \in \dom(F.\omap)$, the stack search in $\sigma$ returns $(F, n)$. By \inv{L1},
    $oid$ is in no region of $H$, so the heap search fails and $\loc{\sigma}{r} = \Stk(n)$. Again by
    \inv{L1}, $oid$ is in no frame of $S$, since $\oids(S :: F) = \oids(S) \cat \dom(F.\omap)$ is
    duplicate-free. So in $\sigma'$ both searches fail, and $\loc{\sigma'}{r} = \noloc$. \qedhere
\end{itemize}
\end{proof}

\begin{corollary}
\label{cor:exit-loc}
For every reference $r$:
\begin{enumerate}[nosep,label=(\alph*)]
  \item $\loc{\sigma}{r} = \Rgn(rid)$ iff $\loc{\sigma'}{r} = \Rgn(rid)$;
    \leanpart{exit\_corollary\_1}
  \item if $\loc{\sigma'}{r} = \Stk(i)$ then $\loc{\sigma}{r} = \Stk(i)$.
    \leanpart{exit\_corollary\_3}
\end{enumerate}
\end{corollary}
\begin{proof}
In the second case of Lemma~\ref{lem:exit-loc}, $\loc{\sigma'}{r}$ is undefined and
$\loc{\sigma}{r}$ is a stack location, so neither side of (a) holds and the premise of (b) is false.
In the first case both statements are trivial.
\end{proof}

\begin{lemma}[No new references]
\label{lem:exit-refs}
\lean{exit\_refs\_mem}
Every element of $\refs(\sigma')$ is an element of $\refs(\sigma)$. More precisely,
$\refs(S) \sublist \refs(S :: F)$ and $\refs(H') \perm \refs(H)$.
\end{lemma}
\begin{proof}
$\refs(S :: F) = \refs(S) \cat \refs(F)$, and $\refs(H') \perm \refs(H)$ is
Lemma~\ref{lem:status-update}(b).
\end{proof}

\subsection{The invariants}

\begin{lemma}[\inv{L1}] \lean{exit\_L1}
$\sigma' \vdash \inv{L1}$.
\end{lemma}
\begin{proof}
$\oids(\sigma') = \oids(S) \cat \oids(H') \perm \oids(S) \cat \oids(H)$ by
Lemma~\ref{lem:status-update}(b), and this is a sublist of
$\oids(S) \cat \dom(F.\omap) \cat \oids(H) = \oids(\sigma)$, which is duplicate-free by \inv{L1}.
\end{proof}

\begin{lemma}[\inv{L2}] \lean{exit\_L2}
$\sigma' \vdash \inv{L2}$.
\end{lemma}
\begin{proof}
Let $G \in S$. By \inv{S1} for $\sigma$, the region ids of $S :: F$ are pairwise distinct, so
$G.\rid \ne F.\rid$. Hence $H'(G.\rid) = H(G.\rid)$, which by \inv{L2} for $\sigma$ (as
$G \in S :: F$) is an $\Open$ region.
\end{proof}

\begin{lemma}[\inv{H1}] \lean{exit\_H1}
$\sigma' \vdash \inv{H1}$.
\end{lemma}
\begin{proof}
A region of $H'$ is either $R'$ or a region of $H$ stored at some key other than $F.\rid$. In the
second case \inv{H1} for $\sigma$ applies directly. In the first, $R'$ has the same $\bridge$ and
$\omap$ as $R$, and $R$ satisfies \inv{H1} since it is a region of $H$.
\end{proof}

\begin{lemma}[\inv{H2}] \lean{exit\_H2}
$\sigma' \vdash \inv{H2}$.
\end{lemma}
\begin{proof}
Fix $rid$ and write $\#\ell$ for the number of occurrences of $\RId(rid)$ in $\ell$. By
Lemma~\ref{lem:exit-refs}, $\#\refs(S) \le \#\refs(S :: F)$ and $\#\refs(H') = \#\refs(H)$, so
$\#\refs(S) + \#\refs(H') \le \#\refs(S :: F) + \#\refs(H) \le 1$ by \inv{H2} for $\sigma$.
\end{proof}

\begin{lemma}[\inv{H3}] \lean{exit\_H3}
$\sigma' \vdash \inv{H3}$.
\end{lemma}
\begin{proof}
Let $H'(rid) = Q'$ and $\OId(oid) \in \refs(Q')$. By Lemma~\ref{lem:status-update}(c),
$H(rid) = Q$ for some $Q$ with the same $\omap$ as $Q'$, and $\refs$ of a region depends only on its
$\omap$, so $\OId(oid) \in \refs(Q)$. By \inv{H3} for $\sigma$,
$\loc{\sigma}{\OId(oid)} = \Rgn(rid)$, and by Corollary~\ref{cor:exit-loc}(a),
$\loc{\sigma'}{\OId(oid)} = \Rgn(rid)$.
\end{proof}

\begin{lemma}[\inv{S1}] \lean{exit\_S1}
$\sigma' \vdash \inv{S1}$.
\end{lemma}
\begin{proof}
$[G.\rid \mid G \in S]$ is a sublist of $[G.\rid \mid G \in S :: F]$, which is duplicate-free by
\inv{S1} for $\sigma$.
\end{proof}

\begin{lemma}[\inv{S2}] \lean{exit\_S2}
$\sigma' \vdash \inv{S2}$.
\end{lemma}
\begin{proof}
Let $(G, i) \in S$, $\OId(oid) \in \refs(G)$ and $\loc{\sigma'}{\OId(oid)} = \Stk(j)$. The pair
$(G,i)$ is also a frame of $\sigma$, at the same index. By Corollary~\ref{cor:exit-loc}(b),
$\loc{\sigma}{\OId(oid)} = \Stk(j)$, so \inv{S2} for $\sigma$ gives $j \le i$.
\end{proof}

\begin{lemma}[\inv{S3}] \lean{exit\_S3}
$\sigma' \vdash \inv{S3}$.
\end{lemma}
\begin{proof}
Let $(G, i) \in S$, $\OId(oid) \in \refs(G)$ and $\loc{\sigma'}{\OId(oid)} = \Rgn(rid')$. As in the
proof of \inv{S2}, $(G, i)$ is a frame of $\sigma$, and by Corollary~\ref{cor:exit-loc}(a)
$\loc{\sigma}{\OId(oid)} = \Rgn(rid')$. \inv{S3} for $\sigma$ gives a frame $(G'', k)$ of $\sigma$
with $G''.\rid = rid'$ and $k \le i$. Since $i < n$ we have $k < n$, so $(G'', k)$ is not the popped
frame: $(G'', k) \in S$, which is the required witness in $\sigma'$.
\end{proof}

\begin{lemma}[\inv{HS1}] \lean{exit\_HS1}
$\sigma' \vdash \inv{HS1}$.
\end{lemma}
\begin{proof}
Let $\OId(oid) \in \refs(\sigma')$. By Lemma~\ref{lem:exit-refs} it is in $\refs(\sigma)$, so by
\inv{HS1} for $\sigma$,
\[ oid \in \oids(\sigma) = \oids(S) \cat \dom(F.\omap) \cat \oids(H). \]
If $oid \in \oids(S)$ it is in $\oids(\sigma')$ directly; if $oid \in \oids(H)$ it is in
$\oids(H') \perm \oids(H)$ (Lemma~\ref{lem:status-update}(b)).

It remains to rule out $oid \in \dom(F.\omap)$: the popped frame's objects must not still be
referenced. Suppose $oid \in \dom(F.\omap)$. By Lemma~\ref{lem:loc-l1}(b),
$\loc{\sigma}{\OId(oid)} = \Stk(n)$. Now $\OId(oid)$ is in $\refs(\sigma') = \refs(S) \cat \refs(H')$:
\begin{itemize}[nosep]
  \item If $\OId(oid) \in \refs(G)$ for some $(G, m) \in S$, then $(G, m)$ is a frame of $\sigma$ and
    \inv{S2} for $\sigma$ gives $n \le m$, contradicting $m < n$.
  \item If $\OId(oid) \in \refs(H')$, then by Lemma~\ref{lem:status-update}(b) it is in $\refs(Q)$
    for some $H(rid) = Q$. \inv{H3} for $\sigma$ gives $\loc{\sigma}{\OId(oid)} = \Rgn(rid)$,
    contradicting $\loc{\sigma}{\OId(oid)} = \Stk(n)$. \qedhere
\end{itemize}
\end{proof}

\begin{lemma}[\inv{HS2}] \lean{exit\_HS2}
$\sigma' \vdash \inv{HS2}$.
\end{lemma}
\begin{proof}
If $\RId(rid) \in \refs(\sigma')$ then $\RId(rid) \in \refs(\sigma)$ by Lemma~\ref{lem:exit-refs}, so
$rid \in \dom(H) = \dom(H')$ by \inv{HS2} for $\sigma$ and Lemma~\ref{lem:status-update}(a).
\end{proof}

\begin{theorem} \lean{exit\_valid}
If $\sigma \vdash \inv{ValidConfig}$ and $\sigma \xrightarrow{\texttt{exit}} \sigma'$, then
$\sigma' \vdash \inv{ValidConfig}$.
\end{theorem}
\begin{proof}
The ten lemmas above.
\end{proof}


%======================================================================
\section{\texttt{enter}}
%======================================================================

\begin{mathpar}
  \inferrule* [lab=enter]
    {\mathit{resolveFA}(xf, \langle S, H\rangle) = \RId(rid) \\ H(rid) = R \\ R.\status = \Closed}
    {\langle S, H \rangle \xrightarrow{\texttt{enter}\ xf\ a} \langle S :: (rid, a, \emptyset, \emptyset),\ H[rid \mapsto R'] \rangle}
\end{mathpar}
where $R' = R$ with $\status := \Open$.

%======================================================================
\section{\texttt{makeObjStack}}
%======================================================================

\begin{mathpar}
  \inferrule* [lab=make-obj-stack]
    {o^* = \mathit{freshObjectId}(\langle S :: F, H\rangle)}
    {\langle S :: F, H \rangle \xrightarrow{\texttt{makeObjStack}\ x} \langle S :: F',\ H \rangle}
\end{mathpar}
where $F' = F$ with $\vmap[x \mapsto \OId(o^*)]$ and $\omap[o^* \mapsto \emptyset]$.

%======================================================================
\section{\texttt{makeObjRegion}}
%======================================================================

\begin{mathpar}
  \inferrule* [lab=make-obj-region]
    {H(F.\rid) = R \\ R.\status = \Open \\ o^* = \mathit{freshObjectId}(\langle S :: F, H\rangle)}
    {\langle S :: F, H \rangle \xrightarrow{\texttt{makeObjRegion}\ x} \langle S :: F',\ H[F.\rid \mapsto R'] \rangle}
\end{mathpar}
where $F' = F$ with $\vmap[x \mapsto \OId(o^*)]$ and $R' = R$ with $\omap[o^* \mapsto \emptyset]$.

%======================================================================
\section{\texttt{makeRegion}}
%======================================================================

\begin{mathpar}
  \inferrule* [lab=make-region]
    {o^* = \mathit{freshObjectId}(\langle S :: F, H\rangle) \\ \rho^* = \mathit{freshRegionId}(\langle S :: F, H\rangle)}
    {\langle S :: F, H \rangle \xrightarrow{\texttt{makeRegion}\ x} \langle S :: F',\ H[\rho^* \mapsto R^*] \rangle}
\end{mathpar}
where $F' = F$ with $\vmap[x \mapsto \RId(\rho^*)]$ and
$R^* = (o^*, [o^* \mapsto \emptyset], \Closed)$.

\texttt{makeRegion}~$x$ allocates a new $\Closed$ region, holding only an empty bridge object, and
writes a reference to the region into the variable $x$ of the top frame. Throughout this section fix
$\sigma = \langle S :: F, H\rangle$ with $\sigma \vdash \inv{ValidConfig}$, and
$\sigma' = \langle S :: F', H'\rangle$ with $H' = H[\rho^* \mapsto R^*]$, where $o^*$, $\rho^*$,
$F'$ and $R^*$ are as above. Let $n = |S|$, the index of $F$ in $\sigma$ and of $F'$ in $\sigma'$.
This is exactly what \texttt{makeRegion\_cases} extracts from
$\texttt{makeRegion}\ x\ \sigma = \mathsf{some}\ \sigma'$. Note that $F'$ has the same region id and
the same $\omap$ as $F$.

The proof has three parts: lemmas that describe the frames, heap, stored references, object ids and
locations of $\sigma'$ in terms of $\sigma$
(Lemmas~\ref{lem:makeRegion-frames}--\ref{lem:makeRegion-loc}), a combination of these for the object
references stored in frames (Lemma~\ref{lem:makeRegion-stack-oid}), and the ten invariants.

\subsection{Supporting lemmas}

\begin{lemma}[Frames after \texttt{makeRegion}]
\label{lem:makeRegion-frames}
\begin{enumerate}[nosep,label=(\alph*)]
  \item For every indexed frame $(G', i)$ of $\sigma'$ there is an indexed frame $(G, i)$ of $\sigma$
    such that every element of $\refs(G')$ is in $\refs(G)$ or equals $\RId(\rho^*)$;
    \leanpart{swi\_down}
  \item for every indexed frame $(G, i)$ of $\sigma$ there is an indexed frame $(G', i)$ of $\sigma'$
    with $G'.\rid = G.\rid$. \leanpart{swi\_up}
\end{enumerate}
\end{lemma}
\begin{proof}
By Lemma~\ref{lem:stack-split}(a), the indexed frames of $\sigma$ are those of $S$ followed by
$(F, n)$, and those of $\sigma'$ are those of $S$ followed by $(F', n)$.

(a) If $(G', i)$ is a frame of $S$, take $G = G'$. Otherwise $(G', i) = (F', n)$; take $(F, n)$. By
Lemma~\ref{lem:write-mem}(e), every element of $\refs(F')$ is $\RId(\rho^*)$ or in $\refs(F)$.

(b) If $(G, i)$ is a frame of $S$, take $G' = G$. Otherwise $(G, i) = (F, n)$; take $(F', n)$, which
has the same region id.
\end{proof}

\begin{lemma}[The heap after \texttt{makeRegion}]
\label{lem:makeRegion-heap}
\begin{enumerate}[nosep,label=(\alph*)]
  \item The entries of $H'$ are $(\rho^*, R^*)$ followed by the entries of $H$;
    \leanpart{heap\_entries}
  \item $\refs(H') = \refs(H)$. \leanpart{heap\_refs\_eq}
\end{enumerate}
\leanform{(a) is stated for an insert of any region at $\mathit{freshRegionId}$.}
\end{lemma}
\begin{proof}
(a) By Lemma~\ref{lem:fresh}(a), $\rho^* \notin \dom(H)$, so the insert prepends an entry and removes
none. (b) By (a), $\refs(H') = \refs(R^*) \cat \refs(H)$, and $\refs(R^*)$ is empty because the only
object of $R^*$ has no fields.
\end{proof}

\begin{lemma}[Object ids after \texttt{makeRegion}]
\label{lem:makeRegion-oids}
\lean{objectIds\_perm}
$\oids(\sigma') \perm [o^*] \cat \oids(\sigma)$.
\end{lemma}
\begin{proof}
By Lemma~\ref{lem:stack-split}(c), Lemma~\ref{lem:makeRegion-heap}(a) and
$\oids(F') = \dom(F.\omap) = \oids(F)$,
\begin{align*}
  \oids(\sigma') &= \oids(S) \cat \dom(F.\omap) \cat [o^*] \cat \oids(H) \\
    &\perm [o^*] \cat \oids(S) \cat \dom(F.\omap) \cat \oids(H),
\end{align*}
and the last line is $[o^*] \cat \oids(\sigma)$.
\end{proof}

\begin{lemma}[References after \texttt{makeRegion}]
\label{lem:makeRegion-refs}
\lean{refs\_mem}
Every element of $\refs(\sigma')$ is $\RId(\rho^*)$ or an element of $\refs(\sigma)$.
\end{lemma}
\begin{proof}
Let $r \in \refs(\sigma')$. By Lemma~\ref{lem:stack-split}(d) and
Lemma~\ref{lem:makeRegion-heap}(b), $r$ is in $\refs(S)$, $\refs(F')$ or $\refs(H)$. In the first and
last cases $r \in \refs(\sigma)$. If $r \in \refs(F')$ then by Lemma~\ref{lem:write-mem}(e) either
$r = \RId(\rho^*)$, or $r \in \refs(F)$, and hence $r \in \refs(\sigma)$.
\end{proof}

\begin{lemma}[Locations after \texttt{makeRegion}]
\label{lem:makeRegion-loc}
\lean{loc\_eq}
For every $oid$,
\[ \loc{\sigma'}{\OId(oid)} =
   \begin{cases}
     \Rgn(\rho^*) & \text{if } oid = o^*, \\
     \loc{\sigma}{\OId(oid)} & \text{otherwise.}
   \end{cases} \]
\end{lemma}
\begin{proof}
The stack search of Definition~\ref{def:loc} looks at the top frame first and then at the frames of
$S$; the heap search looks at the entries of $H'$ in the order given by
Lemma~\ref{lem:makeRegion-heap}(a), so at $(\rho^*, R^*)$ first.
\begin{itemize}[nosep]
  \item If $oid = o^*$: by Lemma~\ref{lem:fresh}(b), $o^*$ is not in $\dom(F.\omap)$, which is the
    $\omap$ of $F'$, nor in the $\omap$ of any frame of $S$. So the stack search fails. The heap
    search stops at $(\rho^*, R^*)$, since $o^* \in \dom(R^*.\omap)$. Hence the location is
    $\Rgn(\rho^*)$.
  \item If $oid \ne o^*$: $oid \notin \dom(R^*.\omap) = \{o^*\}$, so the heap search passes over
    $(\rho^*, R^*)$ and continues exactly as the search of $H$. The stack search first looks at
    $(F', n)$, where it succeeds iff it succeeds at $(F, n)$ in $\sigma$, because the two frames have
    the same $\omap$; in that case both searches return index $n$, and otherwise both continue
    identically through $S$. By Lemma~\ref{lem:loc-searches} the two locations agree. \qedhere
\end{itemize}
\end{proof}

\begin{corollary}
\label{cor:makeRegion-loc}
\lean{makeRegion\_corollary\_loc\_eq}
For every $oid \ne o^*$, $\loc{\sigma}{\OId(oid)} = \loc{\sigma'}{\OId(oid)}$.
\leanform{It takes $\mathit{ValidConfig}\ \sigma$ as a hypothesis but does not use it.}
\end{corollary}
\begin{proof}
The second case of Lemma~\ref{lem:makeRegion-loc}.
\end{proof}

\begin{lemma}[Object references in frames]
\label{lem:makeRegion-stack-oid}
\lean{stack\_oid\_down}
Let $(G', i)$ be an indexed frame of $\sigma'$ with $\OId(oid) \in \refs(G')$. Then there is an
indexed frame $(G, i)$ of $\sigma$ with $\OId(oid) \in \refs(G)$, and
$\loc{\sigma'}{\OId(oid)} = \loc{\sigma}{\OId(oid)}$.
\end{lemma}
\begin{proof}
Take $(G, i)$ as in Lemma~\ref{lem:makeRegion-frames}(a). Since $\OId(oid) \ne \RId(\rho^*)$, we have
$\OId(oid) \in \refs(G)$. So $\OId(oid) \in \refs(\sigma)$ by Lemma~\ref{lem:stored}(d), and
$oid \ne o^*$ by Lemma~\ref{lem:fresh}(c). Lemma~\ref{lem:makeRegion-loc} then gives the equality of
locations.
\end{proof}

\subsection{The invariants}

\begin{lemma}[\inv{L1}] \lean{makeRegion\_L1}
$\sigma' \vdash \inv{L1}$.
\end{lemma}
\begin{proof}
$[o^*] \cat \oids(\sigma)$ is duplicate-free, since $\oids(\sigma)$ is by \inv{L1} for $\sigma$ and
$o^* \notin \oids(\sigma)$ by Lemma~\ref{lem:fresh}(a). By Lemma~\ref{lem:makeRegion-oids},
$\oids(\sigma')$ is a permutation of it.
\end{proof}

\begin{lemma}[\inv{L2}] \lean{makeRegion\_L2}
$\sigma' \vdash \inv{L2}$.
\end{lemma}
\begin{proof}
Let $G' \in S :: F'$. Either $G' \in S$, or $G' = F'$, which has the same region id as $F$. In both
cases \inv{L2} for $\sigma$ gives an $\Open$ region $Q$ with $H(G'.\rid) = Q$. Then
$G'.\rid \in \dom(H)$, so $G'.\rid \ne \rho^*$ by Lemma~\ref{lem:fresh}(a), and
$H'(G'.\rid) = H(G'.\rid) = Q$.
\end{proof}

\begin{lemma}[\inv{H1}] \lean{makeRegion\_H1}
$\sigma' \vdash \inv{H1}$.
\end{lemma}
\begin{proof}
By Lemma~\ref{lem:makeRegion-heap}(a), a region of $H'$ is either $R^*$ or a region of $H$. The
bridge object $o^*$ of $R^*$ is in $\dom(R^*.\omap) = \{o^*\}$. For a region of $H$, \inv{H1} for
$\sigma$ applies directly.
\end{proof}

\begin{lemma}[\inv{H2}] \lean{makeRegion\_H2}
$\sigma' \vdash \inv{H2}$.
\end{lemma}
\begin{proof}
Fix $rid$ and write $\#\ell$ for the number of occurrences of $\RId(rid)$ in $\ell$, and
$\delta$ for $1$ if $rid = \rho^*$ and $0$ otherwise. First,
\[ \#\refs(F') \le \#\refs(F) + \delta. \]
Indeed, $F'$ and $F$ have the same $\omap$, and the values of $F.\vmap[x \mapsto \RId(\rho^*)]$ are
$\RId(\rho^*)$ together with the values of $F.\vmap$ with the entry for $x$ (if any) removed, the
latter being a sublist of the values of $F.\vmap$. Next, by Lemma~\ref{lem:stack-split}(c) and
Lemma~\ref{lem:makeRegion-heap}(b), the quantity to bound is
\[ \#\refs(S :: F') + \#\refs(H') = \#\refs(S) + \#\refs(F') + \#\refs(H), \]
while \inv{H2} for $\sigma$ reads
$\#\refs(S) + \#\refs(F) + \#\refs(H) \le 1$.
\begin{itemize}[nosep]
  \item If $rid \ne \rho^*$, then $\delta = 0$ and the bound follows from the two displays.
  \item If $rid = \rho^*$, then $\RId(\rho^*) \notin \refs(\sigma)$: otherwise \inv{HS2} for $\sigma$
    would give $\rho^* \in \dom(H)$, contradicting Lemma~\ref{lem:fresh}(a). So
    $\#\refs(S) = \#\refs(F) = \#\refs(H) = 0$, hence $\#\refs(F') \le 1$ and the sum is at most
    $1$. \qedhere
\end{itemize}
\end{proof}

\begin{lemma}[\inv{H3}] \lean{makeRegion\_H3}
$\sigma' \vdash \inv{H3}$.
\end{lemma}
\begin{proof}
Let $H'(rid) = Q$ and $\OId(oid) \in \refs(Q)$. If $rid = \rho^*$ then $Q = R^*$, which has no
references, so this case is impossible. Hence $rid \ne \rho^*$ and $H(rid) = Q$. By
Lemma~\ref{lem:stored}(b), $\OId(oid) \in \refs(H) \sublist \refs(\sigma)$, so $oid \ne o^*$ by
Lemma~\ref{lem:fresh}(c). Therefore, by Lemma~\ref{lem:makeRegion-loc} and \inv{H3} for $\sigma$,
$\loc{\sigma'}{\OId(oid)} = \loc{\sigma}{\OId(oid)} = \Rgn(rid)$.
\end{proof}

\begin{lemma}[\inv{S1}] \lean{makeRegion\_S1}
$\sigma' \vdash \inv{S1}$.
\end{lemma}
\begin{proof}
$F'.\rid = F.\rid$, so $[G.\rid \mid G \in S :: F'] = [G.\rid \mid G \in S :: F]$, which is
duplicate-free by \inv{S1} for $\sigma$.
\end{proof}

\begin{lemma}[\inv{S2}] \lean{makeRegion\_S2}
$\sigma' \vdash \inv{S2}$.
\end{lemma}
\begin{proof}
Let $(G', i)$ be a frame of $\sigma'$, $\OId(oid) \in \refs(G')$ and
$\loc{\sigma'}{\OId(oid)} = \Stk(j)$. Lemma~\ref{lem:makeRegion-stack-oid} gives a frame $(G, i)$ of
$\sigma$ with $\OId(oid) \in \refs(G)$ and $\loc{\sigma}{\OId(oid)} = \Stk(j)$. By \inv{S2} for
$\sigma$, $j \le i$.
\end{proof}

\begin{lemma}[\inv{S3}] \lean{makeRegion\_S3}
$\sigma' \vdash \inv{S3}$.
\end{lemma}
\begin{proof}
Let $(G', i)$ be a frame of $\sigma'$, $\OId(oid) \in \refs(G')$ and
$\loc{\sigma'}{\OId(oid)} = \Rgn(rid')$. Lemma~\ref{lem:makeRegion-stack-oid} gives a frame $(G, i)$
of $\sigma$ with $\OId(oid) \in \refs(G)$ and $\loc{\sigma}{\OId(oid)} = \Rgn(rid')$. \inv{S3} for
$\sigma$ gives a frame $(G'', k)$ of $\sigma$ with $G''.\rid = rid'$ and $k \le i$. By
Lemma~\ref{lem:makeRegion-frames}(b), $\sigma'$ has a frame with index $k$ and region id $rid'$,
which is the required witness.
\end{proof}

\begin{lemma}[\inv{HS1}] \lean{makeRegion\_HS1}
$\sigma' \vdash \inv{HS1}$.
\end{lemma}
\begin{proof}
Let $\OId(oid) \in \refs(\sigma')$. It is not $\RId(\rho^*)$, so by Lemma~\ref{lem:makeRegion-refs}
it is in $\refs(\sigma)$. By \inv{HS1} for $\sigma$, $oid \in \oids(\sigma)$, and so
$oid \in \oids(\sigma')$ by Lemma~\ref{lem:makeRegion-oids}.
\end{proof}

\begin{lemma}[\inv{HS2}] \lean{makeRegion\_HS2}
$\sigma' \vdash \inv{HS2}$.
\end{lemma}
\begin{proof}
Let $\RId(rid) \in \refs(\sigma')$, and note that $\dom(H') = \dom(H) \cup \{\rho^*\}$. By
Lemma~\ref{lem:makeRegion-refs}, either $\RId(rid) = \RId(\rho^*)$, so $rid = \rho^* \in \dom(H')$,
or $\RId(rid) \in \refs(\sigma)$, so $rid \in \dom(H) \subseteq \dom(H')$ by \inv{HS2} for $\sigma$.
\end{proof}

\begin{theorem} \lean{makeRegion\_valid}
If $\sigma \vdash \inv{ValidConfig}$ and $\sigma \xrightarrow{\texttt{makeRegion}\ x} \sigma'$, then
$\sigma' \vdash \inv{ValidConfig}$.
\end{theorem}
\begin{proof}
The ten lemmas above.
\end{proof}

%======================================================================
\section{\texttt{varAsgn}}
%======================================================================

In both rules $\sigma = \langle S :: F, H\rangle$.
\begin{mathpar}
  \inferrule* [lab=var-asgn-bridge]
    {\mathit{resolveFA}(y.f, \sigma) = \OId(oid) \\ x = F.\mathit{bridgeVar} \\
     \loc{\sigma}{\OId(oid)} = \Rgn(F.\rid) \\ H(F.\rid) = R}
    {\sigma \xrightarrow{\texttt{x := y.f}} \langle S :: F,\ H[F.\rid \mapsto R'] \rangle}
  \\
  \inferrule* [lab=var-asgn-stack]
    {\mathit{resolveFA}(y.f, \sigma) = \OId(oid) \\ x \ne F.\mathit{bridgeVar} \\
     \mathit{resolveV}(x, \sigma) = \noloc}
    {\sigma \xrightarrow{\texttt{x := y.f}} \langle S :: F',\ H \rangle}
\end{mathpar}
where $R' = R$ with $\bridge := oid$, and $F' = F$ with $\vmap[x \mapsto \OId(oid)]$.

%======================================================================
\section{\texttt{fieldAsgn}}
%======================================================================

In both rules $\sigma = \langle S :: F, H\rangle$ and $n = |S|$, the index of $F$.
\begin{mathpar}
  \inferrule* [lab=field-asgn-stack]
    {\mathit{resolveV}(x, \sigma) = \OId(oid) \\ \mathit{resolveV}(y, \sigma) = \OId(oid_y) \\
     \loc{\sigma}{\OId(oid)} = \Stk(n) \\ F.\omap(oid) = o}
    {\sigma \xrightarrow{\texttt{x.f := y}} \langle S :: F',\ H \rangle}
  \\
  \inferrule* [lab=field-asgn-region]
    {\mathit{resolveV}(x, \sigma) = \OId(oid) \\ \mathit{resolveV}(y, \sigma) = \OId(oid_y) \\
     \loc{\sigma}{\OId(oid)} = \Rgn(F.\rid) \\ \loc{\sigma}{\OId(oid_y)} = \Rgn(F.\rid) \\
     H(F.\rid) = R \\ R.\status = \Open \\ R.\omap(oid) = o}
    {\sigma \xrightarrow{\texttt{x.f := y}} \langle S :: F,\ H[F.\rid \mapsto R'] \rangle}
\end{mathpar}
where $F' = F$ with $\omap[oid \mapsto o[f \mapsto \OId(oid_y)]]$, and $R' = R$ with
$\omap[oid \mapsto o[f \mapsto \OId(oid_y)]]$.

%======================================================================
\section{\texttt{swap}}
%======================================================================

In all four rules $\sigma = \langle S :: F, H\rangle$, $n = |S|$ is the index of $F$,
$\mathit{resolveV}(y, \sigma) = \OId(yoid)$ and $\mathit{resolveFA}(y.f, \sigma) = \mathit{yfRef}$, and
both $\loc{\sigma}{\OId(yoid)}$ and $\loc{\sigma}{\mathit{yfRef}}$ are defined.
\begin{mathpar}
  \inferrule* [lab=swap-stack]
    {x \ne F.\mathit{bridgeVar} \\ F.\vmap(x) = \mathit{xRef} \\
     \loc{\sigma}{\OId(yoid)} = \Stk(n) \\ F.\omap(yoid) = o}
    {\sigma \xrightarrow{\texttt{swap x y.f}} \langle S :: F_s,\ H \rangle}
  \\
  \inferrule* [lab=swap-region-object]
    {x \ne F.\mathit{bridgeVar} \\ F.\vmap(x) = \OId(xoid) \\ \loc{\sigma}{\OId(xoid)} = \Rgn(F.\rid) \\
     \loc{\sigma}{\OId(yoid)} = \Rgn(F.\rid) \\ H(F.\rid) = R \\ R.\status = \Open \\ R.\omap(yoid) = o}
    {\sigma \xrightarrow{\texttt{swap x y.f}} \langle S :: F_r,\ H[F.\rid \mapsto R_r(\OId(xoid))] \rangle}
  \\
  \inferrule* [lab=swap-region-region]
    {x \ne F.\mathit{bridgeVar} \\ F.\vmap(x) = \RId(xrid) \\
     \loc{\sigma}{\OId(yoid)} = \Rgn(F.\rid) \\ H(F.\rid) = R \\ R.\status = \Open \\ R.\omap(yoid) = o}
    {\sigma \xrightarrow{\texttt{swap x y.f}} \langle S :: F_r,\ H[F.\rid \mapsto R_r(\RId(xrid))] \rangle}
  \\
  \inferrule* [lab=swap-region-bridge]
    {x = F.\mathit{bridgeVar} \\ \loc{\sigma}{\OId(yoid)} = \Rgn(F.\rid) \\
     \mathit{yfRef} = \OId(yfoid) \\ \loc{\sigma}{\OId(yfoid)} = \Rgn(F.\rid) \\
     H(F.\rid) = R \\ R.\omap(yoid) = o}
    {\sigma \xrightarrow{\texttt{swap x y.f}} \langle S :: F,\ H[F.\rid \mapsto R_b] \rangle}
\end{mathpar}
where
\begin{itemize}[nosep]
  \item $F_s = F$ with $\vmap[x \mapsto \mathit{yfRef}]$ and $\omap[yoid \mapsto o[f \mapsto \mathit{xRef}]]$;
  \item $F_r = F$ with $\vmap[x \mapsto \mathit{yfRef}]$;
  \item $R_r(v) = R$ with $\omap[yoid \mapsto o[f \mapsto v]]$;
  \item $R_b = R$ with $\bridge := yfoid$ and $\omap[yoid \mapsto o[f \mapsto \OId(R.\bridge)]]$.
\end{itemize}

%======================================================================
\section{\texttt{merge}}
%======================================================================

\begin{mathpar}
  \inferrule* [lab=merge]
    {F.\vmap(x) = \RId(rid') \\ H(F.\rid) = R \\ H(rid') = R' \\ R.\status = \Open \\ R'.\status = \Closed}
    {\langle S :: F, H \rangle \xrightarrow{\texttt{merge}\ x}
     \langle S :: F',\ (H \setminus rid')[F.\rid \mapsto R_m] \rangle}
\end{mathpar}
where $F' = F$ with $\vmap[x \mapsto \OId(R'.\bridge)]$, $R_m = R$ with
$\omap := R.\omap \cup R'.\omap$, and $H \setminus rid'$ is $H$ with the key $rid'$ removed.


\end{document}
